\usepackage{amsmath,amssymb,amsthm,mathtools,xcolor}
\usepackage{hyperref}
\usepackage{thmtools,thm-restate}
\usepackage{enumitem,tikz,algorithm,algorithmic}
\allowdisplaybreaks
%
%
%
%
%

%
%
%
%
%
%
%

%
%

%
%
\newcommand{\IP}{\mathsf{IP}}
\newcommand{\PSPACE}{\mathsf{PSPACE}}
\renewcommand{\P}{\mathsf{P}}
\newcommand{\NP}{\mathsf{NP}}
\newcommand{\QP}{\mathsf{QP}}
\newcommand{\poly}{\mathsf{poly}}

\newcommand{\DiDI}{\mathsf{DiDI}}

\newcommand{\F}{\mathbb{F}}
\newcommand{\N}{\mathbb{N}}
\newcommand{\R}{\mathbb{R}}
\newcommand{\Q}{\mathbb{Q}}
\newcommand{\Z}{\mathbb{Z}}
\newcommand{\C}{\mathbb{C}}
\newcommand{\K}{\mathbb{K}}
\newcommand{\ringR}{\mathsf{R}}

\newcommand{\x}{{\boldsymbol x}}
\newcommand{\alphavec}{\boldsymbol \alpha}
\newcommand{\T}{{\boldsymbol T}}
\newcommand{\y}{{\boldsymbol y}}
\newcommand{\z}{{\boldsymbol z}}
\newcommand{\V}{{\boldsymbol v}}
\newcommand{\vecc}{{\boldsymbol c}}
\newcommand{\e}{{\boldsymbol e}}
\newcommand{\idxvec}{{\boldsymbol i}}
\newcommand{\VBP}{\mathsf{VBP}}
\newcommand{\calC}{\mathcal{C}}
\newcommand{\calD}{\mathcal{D}}
\newcommand{\calH}{\mathcal{H}}
\newcommand{\calU}{\mathcal{U}}
\newcommand{\calT}{\mathcal{T}}
\newcommand{\calM}{\mathcal{M}}
\newcommand{\calJ}{\mathcal{J}}
\newcommand{\calF}{\mathcal{F}}
\newcommand{\val}{\mathsf{val}}

\renewcommand{\sp}{\mathsf{sp}}
\newcommand{\supp}{\mathsf{support}}
\newcommand{\rk}{\mathsf{rk}}
\newcommand{\cf}{\mathsf{coef}}
\newcommand{\pdeg}{\mathsf{deg}}
\newcommand{\dlg}{\mathsf{dlog}}
\newcommand{\size}{\mathsf{size}}
\newcommand{\tdeg}{\mathsf{trdeg}}
%
%
%
\newcommand{\PSP}[2]{\Pi^{#1}\Sigma\Pi^{#2}}
%
\newcommand{\SPSP}[3]{\Sigma^{#1}\Pi^{#2}\Sigma\Pi^{#3}}
%

\newcommand{\WSP}[2]{\wedge^{#1}\Sigma\Pi^{#2}}
%
%
%
\newcommand{\SWSP}[3]{\Sigma^{#1}\!\wedge^{#2}\!\Sigma\Pi^{#3}}
%
\newcommand{\PSU}[2]{\Pi^{#1} \Sigma\wedge^{#2}}
\newcommand{\SPSU}[3]{\Sigma^{#1} \Pi^{#2} \Sigma \wedge^{#3}}
\newcommand{\WSU}[1]{\wedge\Sigma \wedge^{#1}}
\newcommand{\SWSU}[2]{\Sigma^{#1}\!\wedge\!\Sigma\!\wedge^{#2}}
\newcommand{\Gen}[2]{\mathsf{Gen}(#1,#2)}
\newcommand{\SPS}[2]{\Sigma^{#1} \Pi^{#2} \Sigma}

\numberwithin{equation}{section}

\declaretheoremstyle[bodyfont=\normalfont \itshape \upshape,qed=\qedsymbol]{noproofstyle}

\declaretheoremstyle[bodyfont=\normalfont \itshape]{thmstyle}

\declaretheoremstyle[
    postheadhook = {\hspace*{\parindent}}]{thmliststyle}


\declaretheorem[numberlike=equation,style=thmstyle]{theorem}
\declaretheorem[numberlike=equation,style=thmliststyle, name=Theorem]{theoremlist}
\declaretheorem[numberlike=equation,style=noproofstyle,name=Theorem]{theoremwp}
\declaretheorem[name=Theorem,style=thmstyle,numbered=no]{theorem*}
%

\declaretheorem[numberlike=equation,style=thmstyle]{lemma}
\declaretheorem[name=Lemma,numbered=no,style=thmstyle]{lemma*}
\declaretheorem[numberlike=equation,style=noproofstyle,name=Lemma]{lemmawp}

\declaretheorem[numberlike=equation,style=thmstyle]{corollary}
\declaretheorem[name=Corollary,numbered=no,style=thmstyle]{corollary*}
\declaretheorem[numberlike=equation,style=noproofstyle,name=Corollary]{corollarywp}

\declaretheorem[numberlike=equation,style=thmstyle]{proposition}
\declaretheorem[name=Proposition,numbered=no,style=thmstyle]{proposition*}
\declaretheorem[numberlike=equation,style=noproofstyle,name=Proposition]{propositionwp}

\declaretheorem[numberlike=equation,style=thmstyle]{claim}
\declaretheorem[name=Claim,numbered=no,style=thmstyle]{claim*}
\declaretheorem[numberlike=equation,style=noproofstyle,name=Claim]{claimwp}

\declaretheorem[numberlike=equation,style=thmstyle]{fact}
\declaretheorem[name=Fact,numbered=no,style=thmstyle]{fact*}

\declaretheorem[numberlike=equation,style=thmstyle]{problem}
\declaretheorem[name=Problem,numbered=no,style=thmstyle]{problem*}

\declaretheorem[numberlike=equation,style=thmstyle]{conjecture}
\declaretheorem[name=Conjecture,numbered=no,style=thmstyle]{conjecture*}

\declaretheorem[numberlike=equation,style=thmstyle]{definition}
\declaretheorem[unnumbered,name=Definition,style=thmstyle]{definition*}

\declaretheoremstyle[]{defstyle}

\declaretheorem[numberlike=equation,style=defstyle]{remark}
\declaretheorem[unnumbered,name=Remark,style=defstyle]{remark*}
\declaretheorem[unnumbered,style=thmliststyle, name=Remark]{remarklist*}

\renewcommand{\subsectionautorefname}{Subsection}
\renewcommand{\sectionautorefname}{Section}
\providecommand{\chapterautorefname}{Chapter}
\renewcommand{\chapterautorefname}{Chapter}

\newenvironment{myproof}[1]
{\vspace{1ex}\noindent{\em Proof.}\hspace{0.5em}\def\myproof@name{#1}}
{\hfill{\tiny \qed\ (\myproof@name)}\vspace{1ex}}
\newenvironment{proof-sketch}{\medskip\noindent{\em Proof Sketch.}}{\qed\bigskip}
\newenvironment{proof-attempt}{\medskip\noindent{\em Proof attempt.}}{\bigskip}
\newenvironment{proofof}[1]{\medskip\noindent\emph{Proof of #1.}}{\qed\bigskip}

\newcommand{\inparen }[1]{\left(#1\right)}             
\newcommand{\inbrace }[1]{\left\{#1\right\}}           
\newcommand{\insquare}[1]{\left[#1\right]}             
\newcommand{\inangle }[1]{\left\langle#1\right\rangle} 
